\section{The triangular cyclic module and dominance}
\label{sec:dominance}

We now fix the generator from the introduction exactly and state the
cyclic support theorem used in the induction. The second task of this
section is to detect every constituent dominated by the staircase,
using a sign twist and an explicit permutation-module quotient.

Following the triangular arrangement in \cite[Section~4]{ikenmeyer2015},
for $m\ge1$ use the position set
\begin{equation}
 B_m=\{(i,j)\in\mathbb Z_{\ge1}^2:i+j\le m+1\},
 \qquad G_m=S_{B_m}.
 \label{dom:triangle}
\end{equation}
Its row and column set partitions are
\[
 R_m=\{R_{m,i}:1\le i\le m\},\qquad
 R_{m,i}=\{(i,j):1\le j\le m+1-i\},
\]
\[
 C_m=\{C_{m,j}:1\le j\le m\},\qquad
 C_{m,j}=\{(i,j):1\le i\le m+1-j\}.
\]
Order each row by increasing $j$ and each column by increasing $i$.
With $E_m=\C^m$ and its ordered basis $e_1,\ldots,e_m$, form the
alternating block tensors of \Cref{lem:polytabloid} and set
\begin{equation}
 w_m=v_{R_m}\otimes v_{C_m},\qquad
 W_m=\C[G_m]w_m
 \ \subseteq\ E_m^{\otimes B_m}\otimes E_m^{\otimes B_m}.
 \label{dom:cyclic-module}
\end{equation}
Both set partitions have block sizes $m,m-1,\ldots,1$.
Since $\rho_m^t=\rho_m$, \Cref{lem:polytabloid} gives
\begin{equation}
 V_R:=\C[G_m]v_{R_m}\cong S^{\rho_m},\qquad
 V_C:=\C[G_m]v_{C_m}\cong S^{\rho_m},\qquad
 W_m\subseteq V_R\otimes V_C.
 \label{dom:square-inclusion}
\end{equation}
Here and throughout, the action on a tensor product of two $G_m$-modules
is diagonal.  Thus it suffices to prove the following stronger statement.

\begin{theorem}[Cyclic staircase support]
\label{thm:cyclic-saxl}
For every $m\ge1$ and every partition $\lambda\vdash N_m$, the irreducible
$S^\lambda$ occurs in $W_m$.
\end{theorem}

We prove \Cref{thm:cyclic-saxl} by induction in \Cref{sec:completion}.
The present section establishes the case supplied by dominance.

\begin{proposition}
\label{prop:dominated}
If $m\ge1$, $\lambda\vdash N_m$, and $\rho_m\unrhd\lambda$, then
$S^\lambda$ occurs in $W_m$.
\end{proposition}

\begin{proof}
Fix $m$, abbreviate $G=G_m$, and let $\epsilon$ be the one-dimensional
sign representation, also denoting its character by $\epsilon$.
By Young's rule, conjugation of dominance, and the sign-twist identity
in \Cref{tools:dual-sign}, it is enough to show that
$W_m\otimes\epsilon$ contains every constituent of $M^{\rho_m}$.
We will find the cyclic module generated by
$v_{R_m}\otimes v_{R_m}$ inside an isomorphic image of
$W_m\otimes\epsilon$, then map that module onto $M^{\rho_m}$.
The row subgroup
\[
 H_R=\prod_{i=1}^m S_{R_{m,i}}
\]
acts on $v_{R_m}$ by $\epsilon|_{H_R}$.  Consequently the unnormalized signed
diagonal average is exactly
\begin{equation}
 \sum_{h\in H_R}\epsilon(h)\,h w_m
 =v_{R_m}\otimes u,\qquad
 u=\sum_{h\in H_R}h v_{C_m}.
 \label{dom:signed-average}
\end{equation}
To see that $u\ne0$, let $a_0(i,j)=i$ and consider the corresponding word
$e_{a_0}$.  It is fixed by $H_R$.  Its coefficient in $v_{C_m}$ is $1$:
within each increasingly ordered column it is the identity assignment
$1,2,\ldots,m+1-j$.  The coefficient of $e_{a_0}$ in every $h v_{C_m}$ is
therefore also $1$, since $h^{-1}e_{a_0}=e_{a_0}$.  Its coefficient in $u$ is
$|H_R|=\prod_{r=1}^m r!$, which is nonzero.

We next identify the alternating line in $V_R$.  Frobenius reciprocity,
the induction identity for a sign twist, and \Cref{lem:young-rule} give
\begin{align}
 \dim\Hom_{H_R}(\epsilon|_{H_R},\Res_{H_R}^{G}V_R)
 &=\dim\Hom_G(\Ind_{H_R}^{G}\epsilon|_{H_R},V_R)\notag\\
 &=\dim\Hom_G(M^{\rho_m}\otimes\epsilon,V_R)\notag\\
 &=\dim\Hom_G(M^{\rho_m},V_R\otimes\epsilon)=1.
 \label{dom:alternating-line}
\end{align}
For the last equality, $V_R\otimes\epsilon\cong S^{\rho_m}$ by
self-conjugacy of the staircase, and $S^{\rho_m}$ has multiplicity one in
$M^{\rho_m}$.  As $v_{R_m}$ is nonzero and alternating under $H_R$, it spans
the line in \Cref{dom:alternating-line}.

Twisting the diagonal tensor product can be done in its second factor:
the reassociation
\begin{equation}
 (V_R\otimes V_C)\otimes\epsilon
 \longrightarrow V_R\otimes(V_C\otimes\epsilon),\qquad
 (x\otimes y)\otimes z\longmapsto x\otimes(y\otimes z)
 \label{dom:twist-one-factor}
\end{equation}
is a $G$-isomorphism.  Choose a $G$-isomorphism
$\phi:V_C\otimes\epsilon\longrightarrow V_R$, which exists by
\Cref{dom:square-inclusion} and $\rho_m^t=\rho_m$, and choose $0\ne z\in\epsilon$.
Since $u$ is $H_R$-fixed, $u\otimes z$ is alternating under $H_R$.
It follows from \Cref{dom:alternating-line} that
$\phi(u\otimes z)=c\,v_{R_m}$ for some $c\ne0$.
Tensor \Cref{dom:signed-average} with $z$ and apply the reassociation in
\Cref{dom:twist-one-factor}, followed by $\id\otimes\phi$.
The right-hand side becomes $c(v_{R_m}\otimes v_{R_m})$.
Thus the image of $W_m\otimes\epsilon$ under this isomorphism contains
\begin{equation}
 Z_m=\C[G](v_{R_m}\otimes v_{R_m})
 \ \subseteq\ V_R\otimes V_R.
 \label{dom:repeated-row-cyclic}
\end{equation}

We construct a surjection from $Z_m$ onto $M^{\rho_m}$.
Let $E_{\mathrm{out}}$ have basis $f_1,\ldots,f_{2m-1}$ and define the linear map
\[
 q:E_m\otimes E_m\longrightarrow E_{\mathrm{out}},\qquad
 q(e_a\otimes e_b)=f_{a+b-1}.
\]
Regroup the two input layers by position, using ordinary tensor
reassociation, and apply $q$ at every position.  Then project onto the
span $\mathcal M_\alpha$ of words in which letter $r$ occurs exactly $r$
times for $1\le r\le m$, and no letter greater than $m$ occurs.
Thus the content is $\alpha=(1,2,\ldots,m,0,\ldots,0)$ in the output
alphabet.  Denote the resulting linear map by
\begin{equation}
 Q:E_m^{\otimes B_m}\otimes E_m^{\otimes B_m}
 \longrightarrow\mathcal M_\alpha.
 \label{dom:pair-letter-map}
\end{equation}
Both operations commute with position permutations: the same local map
$q$ is used everywhere, and content is unchanged by $G$.  Hence $Q$ is
$G$-equivariant.

Consider any summand in the expansion of $v_{R_m}\otimes v_{R_m}$.
On a row of length $r$, write its two assignments as
$(a_1,\ldots,a_r)$ and $(b_1,\ldots,b_r)$, each a permutation of
$1,\ldots,r$, and put $c_j=a_j+b_j-1$.  Then
\begin{equation}
 \sum_{j=1}^r c_j=r^2,
 \qquad
 \sum_{j=1}^r c_j^2-r^3
 =\sum_{j=1}^r(c_j-r)^2\ge0.
 \label{dom:square-minimum}
\end{equation}
For a term retained by the content projection, the sum of squared output
letters over the whole triangle is $\sum_{r=1}^m r^3$.
Summing \Cref{dom:square-minimum} over the rows forces equality on every
row.  Thus $c_j=r$ at every position in the row of length $r$.
There is only one possible retained output word, namely
\[
 t_m=\bigotimes_{(i,j)\in B_m}f_{m+1-i}.
\]
It remains to compute its coefficient.  On a row of length $r$, each of
the $r!$ assignments $a_j=\pi(j)$ has the unique compatible assignment
$b_j=r+1-\pi(j)$.  The latter is the permutation $\omega_r\pi$, where
$\omega_r(j)=r+1-j$, so the product of the two alternating signs is
\[
 \sgn(\pi)\sgn(\omega_r\pi)
 =\sgn(\omega_r)=(-1)^{\binom r2}.
\]
The row contributions are independent.  In particular, there is no
cancellation, and the exact image is
\begin{equation}
 Q(v_{R_m}\otimes v_{R_m})
 =\left((-1)^{\sum_{r=1}^m\binom r2}\prod_{r=1}^m r!\right)t_m\ne0.
 \label{dom:noncancellation}
\end{equation}
The stabilizer of $t_m$ is exactly $H_R$, since its equal-letter sets are
the rows of $B_m$.  Its orbit is a basis of $\mathcal M_\alpha$, and hence
\begin{equation}
 Q(Z_m)=\C[G]t_m=\mathcal M_\alpha
 \cong\Ind_{H_R}^{G}\mathbf1\cong M^{\rho_m}.
 \label{dom:permutation-quotient}
\end{equation}
This also explains why the increasing content $(1,2,\ldots,m)$ gives the
usual module indexed by the decreasing partition $\rho_m$: reordering
the content parts conjugates the standard Young subgroup.

By \Cref{lem:young-rule}, every $S^\tau$ with $\tau\unrhd\rho_m$ occurs in
the quotient in \Cref{dom:permutation-quotient}.  Semisimplicity and
\Cref{dom:repeated-row-cyclic} imply that it occurs in $W_m\otimes\epsilon$.
If $\rho_m\unrhd\lambda$, reversal of dominance under conjugation gives
$\lambda^t\unrhd\rho_m^t=\rho_m$.  Taking $\tau=\lambda^t$ and twisting
back by $\epsilon$ proves that $S^\lambda$ occurs in $W_m$.
\end{proof}
